% Above: the trade every contact detector makes. Below: the one honest
% comparison available on this bench.
\begin{tikzpicture}[font=\sffamily\scriptsize, x=1cm, y=1cm]

% ================================================================ the trade
\node[chlabel, anchor=west] at (0,4.3) {threshold against response time, with the false-alarm count beside it};

\draw[taxis] (0,0.6) -- (9.4,0.6);
\draw[taxis] (0,0.6) -- (0,3.9);
\node[font=\sffamily\tiny, anchor=east, rotate=90] at (-0.35,2.3) {response time};
\node[tlabel, anchor=north] at (4.7,0.45) {threshold, newton metres};

% response time rises with threshold
\draw[signal, line width=1pt] (0.4,0.85) -- (1.8,1.25) -- (3.4,1.75) -- (5.6,2.5) -- (8.6,3.6);
\foreach \x/\y/\thr/\fa in {0.4/0.85/{0.02}/{1841}, 1.8/1.25/{0.05}/{37},
                            3.4/1.75/{0.10}/{2}, 5.6/2.5/{0.20}/{0}, 8.6/3.6/{0.50}/{0}}{
  \node[sample] at (\x,\y) {};
  \node[font=\sffamily\tiny, anchor=north] at (\x,\y-0.12) {\thr};
  \node[font=\sffamily\tiny, anchor=south, text=red!60!black] at (\x,\y+0.1) {\fa};
}
\node[tlabel, anchor=west, text=red!60!black] at (9.6,3.6) {false alarms in one hour, above each point};
\draw[deadline] (0,3.05) -- (9.4,3.05);
\node[tlabel, anchor=west] at (9.6,3.05) {the 20 ms budget};
\draw[faultedge] (5.6,2.5) -- (5.6,0.7);
\node[tlabel, anchor=west] at (5.75,1.0) {chosen: the first setting with no false alarms};

\node[umlnote, text width=50mm, anchor=north west] at (0,0.1)
  {Every number on this axis is modelled, because the torque step that produces
   the response is injected into the plant. What is real is the shape of the
   trade: a lower threshold always responds sooner and always raises more false
   alarms, and no setting avoids both.};

% ================================================================ the tap
% Shifted down as a block: the note above it is six lines tall and the impulse
% label rises above its own trace.
\begin{scope}[yshift=-13mm]
\node[chlabel, anchor=west] at (0,-2.6) {and the one measurement in the chapter: a tap on the bench};

\node[signame] at (-0.3,-3.4) {inertial unit};
\draw[busline] (0,-3.4) -- (7.0,-3.4);
\draw[signal, line width=1pt, draw=blue!55!black]
  (0,-3.4) -- (2.0,-3.4) -- (2.1,-2.85) -- (2.25,-3.75) -- (2.45,-3.3)
  -- (2.7,-3.45) -- (3.0,-3.4) -- (7.0,-3.4);
\node[sample] at (2.0,-3.4) {};
\node[tlabel, anchor=west] at (3.15,-3.05) {a real impulse, from a finger};
\node[tlabel, anchor=west] at (7.2,-3.4) {detected 2.1 ms later, median of fifty taps};

\node[signame] at (-0.3,-4.6) {the observer};
\draw[busline] (0,-4.6) -- (7.0,-4.6);
\draw[signal, line width=1pt] (0,-4.6) -- (7.0,-4.6);
\node[tlabel, anchor=west] at (7.2,-4.6) {no response, and that is correct};

\node[umlnote, text width=56mm, anchor=north west] at (0,-5.3)
  {The two detectors answer different questions. The inertial unit sees the
   bench move, which really happened. The observer sees a torque the model
   cannot explain, and the model knows nothing about a finger, so there is
   nothing for it to fail to explain. Reporting the non-response is more useful
   than omitting this row.};
\node[umlnote, text width=52mm, anchor=north west] at (6.4,-5.3)
  {This is the only detection in the chapter that involves no model at all, and
   it is included for exactly that reason: one real signal, from real silicon,
   clearly separated from everything above it.};
\end{scope}
\end{tikzpicture}
